\cvsection{Expériences professionnelles}

\begin{cventries}
  \cventry
    {Stagiaire ingénieur de recherche en IA}
    {INRIA}
    {Bordeaux, France}
    {Mars 2026 -- Septembre 2026}
    {
      \begin{cvitems}
        \item {Utilisation de ressources HPC avec Slurm pour réaliser des expérimentations et comparer les performances de plusieurs modèles de langage.}
        \item {Installation et configuration de LLM, puis utilisation de vLLM pour lancer et servir les modèles sur GPU NVIDIA.}
        \item {Évaluation de Qwen, Gemma et GLM et de plusieurs quantifications pour identifier les compromis entre qualité et performances selon l’environnement d’exécution.}
        \item {Gestion d’une application web conteneurisée utilisant des LLM installés localement sur un Mac Studio, avec installation et configuration des modèles.}
        \item {Conteneurisation de composants avec Docker et mise en place d’un pipeline CI/CD sur GitHub pour automatiser leur intégration.}
        \item {Développement d’un outil Python de collecte de métadonnées scientifiques en ligne et transmission des résultats de tâches cron par webhooks.}
        \item {Développement de services Python/FastAPI, d’API REST documentées avec OpenAPI et de la persistance MongoDB.}
        \item {Conception en autonomie, avec le suivi de mes encadrants, d’un agent IA multi-étapes avec orchestration d’outils et gestion de la mémoire et du contexte.}
        \item {Définition du protocole expérimental, mesure de la latence, du TTFT et de la qualité des réponses, puis rédaction d’un article scientifique.}
        \item {Évaluation d’OpenClaw dans un environnement Docker isolé et analyse des risques liés à l’accès autonome aux outils et aux fichiers.}
      \end{cvitems}
    }

  \cventry
    {Stagiaire de recherche en apprentissage par renforcement}
    {Université Karunya}
    {Coimbatore, Inde}
    {Juin 2025 -- Septembre 2025}
    {
      \begin{cvitems}
        \item {Développement d’un modèle d’apprentissage par renforcement et de son environnement d’entraînement pour la navigation autonome d’un robot hospitalier.}
        \item {Déploiement et optimisation du logiciel et du modèle sur un Raspberry Pi afin de le tester en conditions réelles.}
      \end{cvitems}
    }
\end{cventries}
